\chapter{「品質評価」タブの設定}\label{chap:quality}

ここからは、右の設定パネルのタブごとに、設定の意味と数値の決め方を説明します。
各設定の枠の右上に、\textbf{既定値（初期値）}と\textbf{範囲・選択肢}を書いています。\textbf{\color{lsblue}目安}は、迷ったときの出発点です。

\begin{point}[まずは既定値で]
  既定値は、多くの惑星・月の動画でそのままうまくいくように選んであります。最初は変えずに処理し、結果を見てから1つずつ変えてください。
  一度に多くを変えると、何が効いたのか分からなくなります。
\end{point}

品質評価タブの設定は、\textbf{ファイルの読み方}と、\textbf{品質の測り方}を決めます。
ここを変えると、品質評価からやり直しになります。

% ------------------------------------------------------------------------------
\section{各フレームの品質評価}

\sidebyside[0.5]{insp_quality}{\mk[e]{metric}{1}\mk[e]{byteorder}{2}\mk[e]{depth}{3}}{%
  \begin{marklist}
    \item 品質指標
    \item バイトオーダー（SER のみ）
    \item ビット深度の解釈
  \end{marklist}
  ふつうは3つとも既定のままでかまいません。\badge{2}\,\badge{3} は、画像がおかしく見えるときにだけ変えます。}

\begin{setting}{\badge{1} 品質指標}{勾配エネルギー}{勾配エネルギー／周波数帯パワー比}
  フレームの「写りのよさ」を何で測るかを選びます。
  \begin{itemize}[itemsep=0pt, topsep=2pt]
    \item \uil{勾配エネルギー}：明るさの変化（ふち・模様）がどれだけくっきりしているかで測ります。軽くぼかしてから測るので、細かいノイズには反応しにくくなっています。
    \item \uil{周波数帯パワー比}：画像を周波数（模様の細かさ）に分け、中くらい〜細かい模様の強さの割合で測ります。
  \end{itemize}
  \meyasu{まずは \uil{勾配エネルギー}。品質グラフの順位が、自分の目で見た「よいフレーム」と合わないときに \uil{周波数帯パワー比} も試してください。}
\end{setting}

\begin{setting}{\badge{2} バイトオーダー（SERのみ）}{自動判定}{自動判定／little を使う／big を使う}
  16bit の SER ファイルで、2バイトの数値をどちらの順で読むかです。撮影ソフトによってはヘッダの記録が実際と逆のことがあり、
  \LS は中身を見て自動で判定します。
  \meyasu{画像が砂嵐のようにざらざらに壊れて見えるときだけ、\uil{little を使う}／\uil{big を使う} を切り替えてください。
  画面上部の警告の帯のボタンからも切り替えられます（第\ref{sec:byteorder}節）。}
\end{setting}

\begin{setting}{\badge{3} ビット深度の解釈}{ヘッダに従う}{ヘッダに従う／12bit として扱う／14bit として扱う}
  12bit や 14bit のカメラの画像が、16bit の入れ物に入っていることがあります。そのまま 16bit として読むと、画像がとても暗くなります。
  \meyasu{画像が極端に暗く、警告の帯に「12bitデータが16bitコンテナに入っている可能性があります」と出たら、\uil{12bit として扱う} を選びます。}
\end{setting}

% ------------------------------------------------------------------------------
\section{処理範囲}\label{sec:roi}

大きなセンサーで撮り、画面の一部にだけ対象が写っているときに、\textbf{その周りだけ}を処理する機能です。
処理が速くなり、メモリも少なくて済みます。

\sidebyside[0.5]{insp_roi}{\mk[e]{draw}{1}\mk[e]{label}{2}\mk[e]{clear}{3}}{%
  \begin{marklist}
    \item \uil{フレームの上で処理範囲を描く}：オンにして、プレビューの上でドラッグして四角を描きます。
    \item いまの範囲（大きさと位置）が出ます。
    \item \ui{全体に戻す}：範囲を消して、画像全体を処理します。
  \end{marklist}}

使い方と画面の例は第\ref{sec:roi-howto}節で説明します。
範囲を変えると、品質評価からやり直しになります。範囲はファイルごとの設定なので、プリセットには保存されません。

% ------------------------------------------------------------------------------
\section{入力の前処理（範囲・色・補正）}

\sidebyside[0.5]{insp_input}{\mk[e]{range}{1}\mk[e]{bayer}{2}\mk[e]{debayer}{3}\mk[e]{dark}{4}\mk[e]{flat}{5}\mk[e]{clearcal}{6}}{%
  \begin{marklist}
    \item 使うフレームの範囲
    \item 色の配列（Bayer）
    \item デバイヤーの方式
    \item ダーク
    \item フラット
    \item 補正をやめる
  \end{marklist}}

\begin{setting}{\badge{1} 使うフレームの範囲}{空欄（全体）}{1〜最後のフレーム番号}
  動画の一部だけを使います。左に最初、右に最後のフレーム番号を入れます。空欄なら最初から最後まで使います。
  \meyasu{途中で雲が通った、ピントを直した、など、品質グラフの\uil{時系列}で悪い区間がはっきり分かるときに、その区間を外します。
  惑星は自転するので、長すぎる動画（木星なら2〜3分を大きく超えるもの）は範囲を絞ると模様の流れを抑えられます。}
\end{setting}

\begin{setting}{\badge{2} 色の配列（Bayer）}{ヘッダに従う}{ヘッダに従う／モノクロとして扱う／RGGB／GRBG／GBRG／BGGR}
  カラーカメラのセンサーの色の並び方です。ふつうはファイルに記録されているので、自動で正しく読まれます。
  \meyasu{色がおかしい（全体が緑・紫になる、細かい格子模様が出る）ときは、ヘッダが間違っている可能性があります。
  4種類の並びを順に試し、自然な色になるものを選びます。モノクロのカメラなのに色付きで読まれてしまうときは \uil{モノクロとして扱う}。}
\end{setting}

\begin{setting}{\badge{3} デバイヤーの方式}{標準（bilinear）}{標準（bilinear）／高品質（Malvar-He-Cutler）}
  カラーの画素から3色をそろえる計算の方法です。
  \meyasu{\uil{高品質} のほうが細かい模様の色のにじみが少なく、少しシャープです。計算は少し重くなります。ふだんは \uil{標準} で十分です。}
\end{setting}

\begin{setting}{\badge{4}\,\badge{5} ダーク・フラット}{なし}{動画・静止画・静止画のフォルダ}
  \begin{itemize}[itemsep=0pt, topsep=2pt]
    \item \textbf{ダーク}：レンズにふたをして、本番と同じ露出・ゲイン・温度で撮った画像です。センサーの熱かぶりや、いつも明るい画素（ホットピクセル）を引き算で取り除きます。
    \item \textbf{フラット}：明るさが一様な面（薄曇りの空や、白い布ごしの光）を撮った画像です。周辺が暗くなる現象（周辺減光）や、センサーのホコリの影を割り算で補正します。
  \end{itemize}
  \ui{ダークを選ぶ…}\ui{フラットを選ぶ…} で、動画・静止画・静止画のフォルダを選びます。全フレームを平均して「マスター」を作り、デバイヤーの前の画素に掛けます。
  \meyasu{惑星の短い露出では、なくても大きな問題にならないことが多いです。月・太陽の広い視野で周辺減光やホコリの影が気になるときはフラットを、
  長めの露出や高いゲインでホットピクセルが目立つときはダークを使います。}
\end{setting}

\begin{setting}{\badge{6} 補正をやめる}{—}{—}
  選んだダーク・フラットを外します。
\end{setting}

% ------------------------------------------------------------------------------
\section{詳細設定（追跡失敗の判定）}

品質評価のときに、\textbf{使えないフレーム}（惑星が画面の外へ出た、雲で消えた、大きくぶれた、など）を自動で除外する基準です。

\sidebyside[0.5]{insp_qadv}{\mk[w]{outlier}{1}\mk[w]{similarity}{2}\mk[w]{maxshift}{3}}{%
  \begin{marklist}
    \item 外れ値判定の厳しさ（σ）
    \item 参照との類似度の下限
    \item 許容する位置ずれ（px）
  \end{marklist}
  除外したフレームは、アライメントのあとに警告の帯で枚数が出て、\ui{詳細…} で一覧と理由を見られます。}

\begin{setting}{\badge{1} 外れ値判定の厳しさ（σ）}{6.0}{正の数}
  お手本（参照）との似ている度合い（類似度）が、ほかのフレームと比べて飛び抜けて低いフレームを除外する基準です。全フレームの類似度の真ん中の値（中央値）から、ばらつきの何倍（σ、シグマ）離れたら除外するかを決めます。
  \meyasu{小さくするほど厳しく除外します。崩れたフレームが混ざって結果が乱れるときは 4〜5 に下げます。除外されすぎるときは 8〜10 に上げます。}
\end{setting}

\begin{setting}{\badge{2} 参照との類似度の下限}{0.5}{0〜1}
  お手本（参照）と比べた似ている度合い（1 が完全に同じ）がこの値より低いフレームを除外します。
  \meyasu{雲で薄くなったフレームなどを除きたいときは 0.6〜0.7 に上げます。ゆらぎが大きい日にフレームが除外されすぎるときは 0.3〜0.4 に下げます。}
\end{setting}

\begin{setting}{\badge{3} 許容する位置ずれ（px）}{空欄（自動：画像の短い辺の 1/4）}{画素数}
  お手本からこれ以上ずれたフレームを除外します。
  \meyasu{架台の追尾が悪く、惑星が画面の中で大きく動いてしまった動画で、フレームが除外されすぎるときは大きくします。}
\end{setting}
